\section{The remaining premise of the marked theorem}

% BEGIN A71M1
\paragraph{A71M1. The prime tests for arbitrary chains.}
First take a normalized instance of the original $\mathrm{Core}$.
Its assumptions include $\ell\ne2$, $T.\mathrm{ProL}$, adaptedness,
the coefficient generation and section, and the given patched
points carrying their bounds and ordered eigenvalues. At a vertex
let $W$ be the limit selected by
LIM18 for this $T$. R70P1 places the lifted point on the
prescribed prime and orders the auxiliary roots there. Apply
L71C1--5 on this same $W$, with the bounds $B1,B2$ and the
auxiliary eigenvalues of the vertex. We obtain the three
required Laurent-coefficient Selmer statements. U71R6 gives
the polarized universal matrices on this vertex's completed
ring. The native coefficient DVR and residual identification
are the ones described in R70P5. With this data the local
context has the ordered characters and root relations of
U71R5--6, and the three bounds just proved. P71S6 now gives
the generators, independent rows and relations of its
$\mathrm{Pres}(D_0)$ on exactly this ring.

R70P4--5 prove the local prime containment from those matrices
and that presentation. This covers every two-block vertex.
R70P6 also gives the comparison directly for a full block.
R70P2--3 join successive vertices over the predetermined
connecting primes, so the final containment holds for every
allowed endpoint and the \emph{same} initial Hida ideal.
This proves $\mathrm{Core}$. BCN23 gives
$\mathrm{CoreBC}$ by changing coefficients and introducing the
section. PQL20 then gives $\mathrm{Raw}$ for an arbitrary original
system: it constructs a relative pro-$\ell$ quotient with
the required group property, transports the data there and returns
the containment to the original kernels. Thus $\mathrm{Raw}$ has
its original scope. Finally use H70M1: N69E1--3 give the
normalization, Q70C7 the simultaneous detection, and T70B4
and U70O8 both bounds on the original arithmetic auxiliary
system. N69E4 checks the pointwise residual-kernel condition required
for these arithmetic bounds. Apply the just-proved $\mathrm{Raw}$
to the original auxiliary system. Hence N69E5 gives
$\mathrm{P1core}(J)$ and $\mathrm{P3core}(J)$ for each
input to $\mathrm{PfromAut}$. This implication now needs
no further exceptional-vertex construction.
% END A71M1

% BEGIN A71M2
\paragraph{A71M2. From the Hida supplier to the whole marked conclusion.}
Suppose $\mathrm{HidaUAlt}$ for the original arithmetic
setups. Given its ideal $J$ and automorphy system, A71M1
gives both prime tests. G71H4 supplies $\mathrm{CohP}$,
and KX3 supplies $\mathrm{LocCFTAll}$. HIA3 now yields
$\mathrm{ModUAlt}$; the same proof only calls the supplier at the
setups occurring there. G71H3 gives the cohomological input for
AFE16 on every marked host. L69D3 and D70W5 give the linked
conclusions and the required descent. Applying the marked comparison
of M65A4 as in H70M2 yields
$\mathrm{ThmS6}_s$. The mark is fixed before all universal rings
and all type-I points, with both conclusions on that same mark.
In particular this proves the implication
\[
          \mathrm{HidaUAlt}\ \Longrightarrow\ \mathrm{ThmS6}_s
\]
on the classical foundations specified in these chapters.
The result already includes the general patching and the
cohomological host.

HIA2 constructs $\mathrm{HidaUAlt}$ from the source
$H'=\mathrm{Af.HypothesesH}'$. Of this package F71A1--3 now give
the definite quaternion algebra, the weak unitary descent with
its specified inert-place vectors, and the entire cohomology
entry of the rank-general literature hypothesis. The latter is
only one of that hypothesis's four entries. In particular the
other local representation theory, automorphic and CM entries
and the other fields of $H'$ still require verification in
their precise formulations, or a direct construction of the
needed Hida ideals. Neither alternative is a premise of the
original marked contract. Its status therefore remains open.
% END A71M2
